\section{Primitive middle-degree targets for \texorpdfstring{\Fiii{}}{(F3')}}
\label{v7:sec:primitive-middle}

\begin{theorem}[Primitive middle-degree reduction]
\label{v7:thm:primitive-middle}
The general statement \Fiii{} is equivalent to its restriction to
primitive rational Hodge classes in the middle cohomology of smooth
projective varieties of even dimension. It is also equivalent to the
same restriction on smooth projective rational varieties.
\end{theorem}

\begin{proof}
Let \(Y\) be smooth projective of dimension \(d\), \(L\) an ample class,
and \(\alpha\in H^{2p}(Y,\QQ)\) a primitive Hodge class with \(2p\le d\).
Set \(r=d-2p\), \(W=Y\times\mathbb P^r\), and let
\(\pi\colon W\to Y\) and \(q\colon W\to\mathbb P^r\) be the projections.
Write \(h\) for the hyperplane class on \(\mathbb P^r\). Define the
algebraic class
\[
 D_r=\sum_{j=0}^{r}(-1)^j\pi^*L^j\,q^*h^{r-j}\ \in\operatorname{CH}^r(W)_{\QQ}.
\]
Here the same letters denote the divisor classes in the Chow ring and
their cohomology classes. The transpose of the graph of \(\pi\),
multiplied by the pullback of \(D_r\) from the \(W\)-factor, is a
correspondence from \(Y\) to \(W\) inducing
\(\alpha\mapsto\Theta=D_r\cup\pi^*\alpha\).

The dimension of \(W\) is \(d+r=2(p+r)\), so \(\Theta\) is a
middle-degree Hodge class. Multiplication by \(\pi^*L+q^*h\) cancels
the adjacent terms of the sum. The end terms vanish because
\(L^{r+1}\alpha=0\) and \(h^{r+1}=0\). The sum of the two pullback
ample classes is ample on \(W\), so \(\Theta\) is primitive for this
polarization. Integration on \(\mathbb P^r\) gives
\(\pi_*\Theta=\alpha\), since \(\int_{\mathbb P^r} h^r=1\) and the
terms with \(j>0\) have lower projective degree.

If \(\Theta=\sum_i(\Gamma_i)_*\beta_i\) comes from abelian varieties,
then
\[
 \alpha=\sum_i(\Gamma_\pi\circ\Gamma_i)_*\beta_i,
\]
where \(\Gamma_\pi\) is the graph of \(\pi\).
The composition is algebraic. This gives \Fiii{} on primitive classes
of every degree below or at the middle. Lefschetz decomposition is
rational and respects Hodge structures with Tate twists. It expresses
each Hodge class, in any degree, as a sum of powers of \(L\) applied
to primitive Hodge classes of degree at most \(d\). Cup product with
an algebraic divisor is induced by an algebraic correspondence.
This proves the first equivalence.

If the primitive assertion is assumed only for rational \(W\), the
same argument first gives \Fiii{} in all degrees on rational \(Y\),
since \(Y\times\mathbb P^r\) is rational. For arbitrary \(Y\), embed it
in \(\mathbb P^N\) with codimension \(c\ge2\). The blow-up \(B\) of
\(\mathbb P^N\) along \(Y\) is smooth, projective and rational.
Let \(j\colon E\hookrightarrow B\) be the exceptional divisor,
\(e\colon E\to Y\) its projection, and
\(\zeta=c_1(\mathcal O_E(1))\). The maps
\[
 a(\gamma)=j_*e^*\gamma,\qquad
 b(\delta)=-e_*\!\left(\zeta^{c-2}j^*\delta\right)
\]
are induced by algebraic correspondences and shift codimension by
\(+1\) and \(-1\), respectively. The self-intersection formula gives
\(j^*j_*e^*\gamma=-\zeta e^*\gamma\), and the projective bundle formula
gives \(e_*(\zeta^{c-1})=1\). Hence \(b\circ a=\mathrm{id}\).
Represent \(a(\gamma)\) by \Fiii{} on \(B\), then compose its
correspondences with \(b\) to represent \(\gamma\) on \(Y\).
The reverse implications are restrictions of the universal statement.
\end{proof}

\begin{figure}[ht]
\centering
\begin{tikzpicture}[>=Stealth, every node/.style={font=\small}]
\node (Y) at (0,0) {$\alpha\in H^{2p}(Y)$};
\node (W) at (7.0,0) {$\Theta\in H^{\dim W}(W)$ primitive};
\draw[->] (Y) to[bend left=20] node[above] {$D_r\pi^*$} (W);
\draw[->] (W) to[bend left=20] node[below] {$\pi_*$} (Y);
\node[align=center] at (3.5,-1.3) {$\pi_*\Theta=\alpha$\quad and\quad $(L+h)\Theta=0$};
\end{tikzpicture}
\caption{An explicit passage to primitive middle degree and back.
The arrows are algebraic correspondences. Neither arrow constructs
the required abelian source for the target class.}
\label{fig:primitive-middle-target}
\end{figure}

The graph and blow-up constructions move a class to a primitive
middle-degree target and recover the original class. The remaining
step is to produce an abelian source and an algebraic correspondence
whose image contains that target. Primitivity alone gives no such
correspondence.
